% CCTT Advantages — Provable Asymptotic and Effort Comparisons
% with Traditional Proof Assistants (F*, Lean 4)

\chapter{Comparative Advantages of CCTT over General-Purpose Proof Assistants}
\label{ch:advantages}

This chapter establishes rigorous comparison theorems between CCTT and
general-purpose proof assistants (Lean~4, F${}^*$, Coq, Agda) on the
specific task of \emph{Python program equivalence verification}.  We
restrict attention to results that are both \emph{mathematically clean}
(formalised with precise definitions and complete proofs) and
\emph{adversarially robust} (we explicitly address the strongest
counter\-arguments against each claim).

\medskip
\textbf{Scope restriction.}  All theorems in this chapter concern the
\emph{OTerm-expressible fragment} of Python — the set of programs that
the CCTT compiler (Chapter~\ref{ch:compilation}) successfully maps to
the OTerm intermediate language.  CCTT makes no claim about programs
outside this fragment, and all comparisons with Lean/F${}^*$ are
restricted to the \emph{same} verification task (input-output
equivalence of two Python functions) rather than general theorem
proving.

%% ═══════════════════════════════════════════════════════════
\section{Formal Setup}
\label{sec:adv-setup}
%% ═══════════════════════════════════════════════════════════

\begin{definition}[The OTerm-expressible fragment]
\label{def:oterm-fragment}
Let $\mathcal{L}_{\mathrm{OT}}$ denote the set of Python functions
$f$ for which the CCTT compiler produces a non-$\bot$ OTerm:
\[
  \mathcal{L}_{\mathrm{OT}}
  \;\coloneqq\;
  \bigl\{\, f \in \mathrm{PyFunc}
  \;\bigm|\;
  \mathrm{compile\_to\_oterm}(\mathit{src}(f)) \neq \bot
  \,\bigr\}.
\]
The OTerm language includes: variables ($\mathsf{OVar}$), literals
($\mathsf{OLit}$), $n$-ary operations ($\mathsf{OOp}$), folds
($\mathsf{OFold}$), case splits ($\mathsf{OCase}$), fixed points
($\mathsf{OFix}$), sequences ($\mathsf{OSeq}$), dictionaries
($\mathsf{ODict}$), abstractions ($\mathsf{OLam}$), mapped
collections ($\mathsf{OMap}$), exception handling ($\mathsf{OCatch}$),
quotient types ($\mathsf{OQuotient}$), and abstract specifications
($\mathsf{OAbstract}$).
\end{definition}

\begin{definition}[Verification task]
\label{def:verification-task}
A \emph{Python equivalence verification task} is a pair
$(f, g) \in \mathrm{PyFunc}^2$ together with the question:
\[
  \forall\, \vec{x} \in \mathrm{dom}(f) \cap \mathrm{dom}(g).\;
  f(\vec{x}) = g(\vec{x}) \;\;?
\]
We write $\mathrm{Equiv}(f, g)$ for the proposition that this holds.
\end{definition}

\begin{definition}[Human annotation effort]
\label{def:annotation-effort}
For a verification system $\mathcal{V}$ and task $(f,g)$, the
\emph{human annotation effort} $\mathrm{HA}_{\mathcal{V}}(f,g)$ is
the number of tokens (in the system's input language) that a human
must write \emph{beyond} the source code of $f$ and $g$ to obtain a
machine-checked verdict.

For CCTT, the input is two Python source strings.  For Lean~4, the
input is two Lean definitions plus a proof term or tactic script.
\end{definition}

\begin{definition}[Normalizer-decidable fragment]
\label{def:normalizer-decidable}
Let $\mathcal{N} \subseteq \mathcal{L}_{\mathrm{OT}}^2$ denote the
set of program pairs $(f,g)$ for which the 7-phase normalizer
produces identical canonical forms:
\[
  \mathcal{N}
  \;\coloneqq\;
  \bigl\{\, (f, g) \in \mathcal{L}_{\mathrm{OT}}^2
  \;\bigm|\;
  \mathrm{normalize}(\mathrm{OT}(f)).\mathrm{canon}()
  =
  \mathrm{normalize}(\mathrm{OT}(g)).\mathrm{canon}()
  \,\bigr\},
\]
where $\mathrm{OT}(f)$ denotes the compiled OTerm (after parameter
renaming).
\end{definition}

\begin{definition}[Path-decidable fragment]
\label{def:path-decidable}
Let $\mathcal{P}_d \supseteq \mathcal{N}$ denote the set of program
pairs for which the bidirectional BFS path search
(Chapter~\ref{ch:path-search}) with depth bound $d$ and axiom set
$\mathcal{A}$ finds a connecting path:
\[
  \mathcal{P}_d
  \;\coloneqq\;
  \bigl\{\, (f, g) \in \mathcal{L}_{\mathrm{OT}}^2
  \;\bigm|\;
  \mathrm{search\_path}(\mathrm{nf}(f), \mathrm{nf}(g),
    \text{depth}=d).\mathrm{found} = \mathrm{True}
  \,\bigr\},
\]
where $\mathrm{nf}(f) \coloneqq \mathrm{normalize}(\mathrm{OT}(f))$.
Note $\mathcal{N} = \mathcal{P}_0$.
\end{definition}


%% ═══════════════════════════════════════════════════════════
\section{Theorem 1: Zero Human Annotation on the Normalizer-Decidable
         Fragment}
\label{sec:thm-annotation}
%% ═══════════════════════════════════════════════════════════

\begin{theorem}[Zero-annotation verification]
\label{thm:zero-annotation}
For every $(f, g) \in \mathcal{P}_d$ (Definition~\ref{def:path-decidable}),
CCTT decides $\mathrm{Equiv}(f,g)$ with
\[
  \mathrm{HA}_{\mathrm{CCTT}}(f, g) = 0.
\]
For the same task in Lean~4, let $n \coloneqq |\mathit{src}(f)| +
|\mathit{src}(g)|$ be the combined source size.  Then
\[
  \mathrm{HA}_{\mathrm{Lean}}(f, g)
  \;\geq\; \Omega(n),
\]
where the lower bound accounts for the \emph{mandatory} Lean
formalisation of both $f$ and $g$.
\end{theorem}

\begin{proof}
\textbf{CCTT side.}
The CCTT pipeline (\texttt{check\_equivalence}) takes two raw Python
source strings as input and returns a verdict.  The user provides
no type annotations, no proof scripts, no lemma statements, and no
intermediate hints.  The pipeline:
\begin{enumerate}[label=(\roman*)]
  \item compiles each source to an OTerm via
    $\mathrm{compile\_to\_oterm}$,
  \item applies the 7-phase normaliser (Defn.~\ref{def:normalizer-decidable}),
  \item if canonical forms match and pass the quality gates
    (no OUnknown, sufficient parameter references, no unresolved
    names), returns $\mathsf{True}$,
  \item otherwise, attempts path search (Defn.~\ref{def:path-decidable})
    with depth bound $d$ and axiom set $\mathcal{A}$;
    returns $\mathsf{True}$ if a path is found.
\end{enumerate}
No step requires human input beyond the two source strings.  Hence
$\mathrm{HA}_{\mathrm{CCTT}}(f,g) = 0$.

\textbf{Lean side.}
To verify $\mathrm{Equiv}(f, g)$ in Lean~4, the user must:
\begin{enumerate}[label=(\alph*)]
  \item \emph{Translate} $f$ and $g$ into Lean definitions.  Since
    there is no general-purpose Python-to-Lean translator that
    preserves semantic equivalence for the OTerm fragment, the user
    must produce Lean definitions that faithfully model each Python
    function.  The translation length is $\Omega(n)$ — each
    syntactic construct in the Python source must map to a
    corresponding Lean construct.
  \item \emph{State the equivalence lemma}: at minimum,
    \texttt{theorem equiv : $\forall$ x, f x = g x}.  This is $O(1)$
    tokens.
  \item \emph{Produce a proof}: either a term-mode proof or a tactic
    script.  Even if the proof is trivially \texttt{by decide} or
    \texttt{by rfl}, steps (a) and (b) are still required.
\end{enumerate}
Step (a) contributes $\Omega(n)$ tokens to the human annotation
effort, giving $\mathrm{HA}_{\mathrm{Lean}}(f,g) \geq \Omega(n)$.
\end{proof}

\begin{remark}[Adversarial analysis of Theorem~\ref{thm:zero-annotation}]
\label{rem:adv-annotation}
\textbf{Objection 1: ``Lean could also take Python source directly.''}
In principle, one could build a Python-to-Lean compiler (cf.\
\cite{greenman2014getting}).  However, no such compiler currently
exists for the OTerm fragment, and building one is itself a
substantial engineering and verification effort.  If such a compiler
existed and were trusted, it would reduce $\mathrm{HA}_{\mathrm{Lean}}$
to the proof effort alone — but the proof effort remains nonzero
for any pair not in Lean's definitional equality kernel.

\textbf{Objection 2: ``Lean's \texttt{decide} tactic gives O(1) proof effort.''}
This is true for decidable propositions over finite types.  However,
$\texttt{decide}$ requires that the proposition be in the
$\texttt{Decidable}$ type class, which requires a decision procedure
in Lean's kernel.  For Python program equivalence over $\PyObj$
(which includes integers, strings, lists, and dictionaries), the
required decision procedure does not exist in Lean's library.  The
user would need to \emph{implement} it — which is precisely the work
that CCTT's normaliser and path search perform.

\textbf{Objection 3: ``The comparison is unfair because CCTT is specialised.''}
This is correct and is a feature, not a bug.  Theorem~\ref{thm:zero-annotation}
does not claim that CCTT is ``better'' than Lean in general.  It
claims that for the \emph{specific} task of Python equivalence
verification on the OTerm fragment, CCTT eliminates the human
annotation cost that Lean requires.  This is analogous to how a
domain-specific SAT solver can outperform a general-purpose theorem
prover on satisfiability problems without being ``better'' overall.

\textbf{What is genuinely new.}
The mathematical content is that the normaliser-plus-path-search
pipeline constitutes a \emph{decision procedure} for the fragment
$\mathcal{P}_d$, whereas Lean's kernel does not contain a decision
procedure for this fragment.  Decision procedures eliminate human
effort; their absence requires it.
\end{remark}


%% ═══════════════════════════════════════════════════════════
\section{Theorem 2: Monomorphic Fiber Decomposition Avoids
         Theory Combination Overhead}
\label{sec:thm-fiber}
%% ═══════════════════════════════════════════════════════════

Before stating the theorem, we establish the complexity landscape.

\begin{definition}[Type fiber decomposition]
\label{def:type-fiber-decomp}
Let $f : \PyObj^m \to \PyObj$ be a compiled function with $m$
parameters.  The \emph{type fibration}
$\pi : \PyObj \to \mathsf{TypeTag}$ (where
$\mathsf{TypeTag} = \{\mathsf{Int}, \mathsf{Bool}, \mathsf{Str},
\mathsf{Pair}, \mathsf{Ref}, \mathsf{None}\}$) induces a covering
\[
  \PyObj^m = \bigcup_{\vec{\tau} \in \mathsf{TypeTag}^m}
  U_{\vec{\tau}},
  \qquad
  U_{\vec{\tau}}
  \coloneqq
  \bigl\{\, \vec{x} \in \PyObj^m
  \;\bigm|\;
  \pi(x_i) = \tau_i \;\forall i
  \,\bigr\}.
\]
The number of fibers is $k \coloneqq |\mathsf{TypeTag}|^m = 6^m$,
but duck-type inference (Chapter~\ref{ch:opacity-zone}) reduces
this to the set of \emph{relevant} fibers $\mathcal{R} \subseteq
\mathsf{TypeTag}^m$ with $|\mathcal{R}| \leq k$.
\end{definition}

\begin{definition}[Monomorphic SMT problem]
\label{def:mono-smt}
A \emph{monomorphic SMT problem over fiber $U_{\vec{\tau}}$} is a
Z3 satisfiability query of the form
\[
  \exists \vec{x} \in U_{\vec{\tau}}.\;
  \den{f}(\vec{x}) \neq \den{g}(\vec{x}),
\]
where the type tag constraints
$\bigwedge_i \mathsf{tag}(x_i) = \tau_i$ restrict each variable to
a single constructor of the $\PyObj$ ADT.

On fiber $\mathsf{Int}^m$: the problem reduces to quantifier-free
linear integer arithmetic (QF\_LIA) when all operations are
arithmetic, which is decidable in $\mathrm{NP}$.

On fiber $\mathsf{Str}^m$: the problem reduces to the quantifier-free
theory of strings (QF\_S), also decidable.

On fiber $\mathsf{Bool}^m$: the problem reduces to propositional
SAT, decidable in $\mathrm{NP}$.
\end{definition}

\begin{definition}[Polymorphic SMT problem]
\label{def:poly-smt}
The \emph{polymorphic (un-decomposed) SMT problem} is the single query
\[
  \exists \vec{x} \in \PyObj^m.\;
  \den{f}(\vec{x}) \neq \den{g}(\vec{x}),
\]
with no type tag constraints.  This query involves the \emph{full}
$\PyObj$ ADT (7 constructors) and potentially all Z3 theory fragments
simultaneously — integers, strings, booleans, and algebraic datatypes
— requiring the Nelson--Oppen theory combination procedure.
\end{definition}

\begin{theorem}[Fiber decomposition reduces theory combination]
\label{thm:fiber-decomp}
Let $(f, g) \in \mathcal{L}_{\mathrm{OT}}^2$ be a verification task
with $m$ parameters, and let $\mathcal{R}$ be the set of relevant
type fibers determined by duck-type inference, with
$|\mathcal{R}| = k$.

\begin{enumerate}[label=(\alph*)]
  \item \textbf{Monomorphic queries are in simpler theory fragments.}
    On each fiber $U_{\vec{\tau}}$ with $\vec{\tau} \in \mathcal{R}$,
    the monomorphic SMT query
    $\exists \vec{x} \in U_{\vec{\tau}}.\; \den{f}(\vec{x}) \neq
    \den{g}(\vec{x})$ lies in a \emph{single-sorted} decidable
    fragment (QF\_LIA for $\mathsf{Int}^m$, QF\_S for
    $\mathsf{Str}^m$, QF\_BV/SAT for $\mathsf{Bool}^m$).  No
    theory combination is required within each fiber query.

  \item \textbf{The polymorphic query requires theory combination.}
    The un-decomposed query over $\PyObj^m$ involves the ADT sort
    (matching on constructors), integer arithmetic, string operations,
    and boolean logic simultaneously.  Z3 must use the Nelson--Oppen
    combination procedure (or DPLL($T_1 + \cdots + T_j$)) to combine
    $j \geq 2$ theory solvers.

  \item \textbf{Composition via $\Hone$.}
    The local verdicts compose into a global verdict via the \v{C}ech
    $\Hone$ computation, which requires $O(k^2)$ overlap checks and
    a GF(2) Gaussian elimination of rank $O(k)$.  This overhead is
    polynomial in $k$ and independent of the SMT problem size.

  \item \textbf{Resulting complexity comparison.}
    Let $T_{\mathrm{mono}}(\vec{\tau}, n)$ be the time for Z3 to solve
    the monomorphic query on fiber $\vec{\tau}$ with $n$ operations in
    the Z3 term.  Let $T_{\mathrm{poly}}(n)$ be the time for the
    polymorphic query.  Then CCTT's total verification time is
    \[
      T_{\mathrm{CCTT}} = O\!\left(
        \sum_{\vec{\tau}\in\mathcal{R}}
          T_{\mathrm{mono}}(\vec{\tau}, n)
        + k^2
      \right),
    \]
    whereas a monolithic approach requires $T_{\mathrm{poly}}(n)$,
    which involves the overhead of Nelson--Oppen combination.
\end{enumerate}
\end{theorem}

\begin{proof}
(a) On fiber $U_{\vec{\tau}}$, each parameter $x_i$ is constrained
to a single $\PyObj$ constructor.  The Z3 term
$\den{f}(\vec{x})$ on this fiber involves only the operations
relevant to that constructor: for $\mathsf{Int}$, these are
$+, -, \times, //, \bmod, <, \leq, >, \geq, =, \neq$, which lie in
QF\_LIA.  The $\PyObj$ ADT constructors become ``transparent'' —
$\mathsf{ival}(\mathsf{IntObj}(n)) = n$ is definitional in Z3.
Hence the query is equivalent to a query in the single-sorted fragment.

Formally: let $\sigma_\tau : U_\tau \hookrightarrow \PyObj$ be the
fibre inclusion.  The pullback $\sigma_\tau^* \den{f}$ replaces every
$\mathsf{tag}$ test with its known value, eliminating constructor
matching.  The resulting term mentions only operations of theory
$T_\tau$, so Z3 can solve it using only the $T_\tau$ solver.

(b) Without type constraints, the Z3 term $\den{f}(\vec{x})$
contains $\mathsf{If}(\mathsf{tag}(x) = \mathsf{Int}, \ldots,
\mathsf{If}(\mathsf{tag}(x) = \mathsf{Str}, \ldots))$ case splits.
This forces Z3 to reason simultaneously about integer arithmetic
(inside the $\mathsf{Int}$ branch), string operations (inside the
$\mathsf{Str}$ branch), and the ADT structure (the case splits
themselves).  The Nelson--Oppen combination procedure propagates
equalities between theory solvers, incurring overhead proportional
to the number of shared variables and the number of theories.

(c) The \v{C}ech complex has:
$C^0 \coloneqq \bigoplus_{\vec{\tau}} \mathbb{F}_2$ with $k$ basis
elements (one per fiber);
$C^1 \coloneqq \bigoplus_{\text{overlaps}} \mathbb{F}_2$ with
$\binom{k}{2}$ potential entries.
The coboundary $\delta^0 : C^0 \to C^1$ is a $\binom{k}{2} \times k$
matrix over $\mathrm{GF}(2)$.  Gaussian elimination runs in
$O(k^3)$ time.  Since $k \leq 6^m$ and $m$ is typically small
(1-3 parameters), this is negligible.

(d) Combining (a)--(c): CCTT solves $k$ independent single-theory
queries plus a polynomial gluing computation; the monolithic approach
solves one multi-theory query.  Each monomorphic query avoids the
Nelson--Oppen overhead that the polymorphic query incurs.
\end{proof}

\begin{remark}[Adversarial analysis of Theorem~\ref{thm:fiber-decomp}]
\label{rem:adv-fiber}
\textbf{Objection 1: ``Z3 already does case splitting internally.''}
Z3's DPLL(T) engine \emph{does} case-split on ADT constructors.
However, this is demand-driven and interleaved with theory propagation
— Z3 does not a priori decompose into independent subproblems.  In
contrast, CCTT's decomposition is \emph{eager} and \emph{external}:
each fiber yields a completely independent Z3 invocation with its
own solver instance and timeout.  The advantages of eager
decomposition are:
\begin{enumerate}[label=(\roman*)]
  \item \emph{Parallelism}: the $k$ fiber queries are independent
    and can be solved in parallel, giving wall-clock time
    $\max_{\vec{\tau}} T_{\mathrm{mono}}(\vec{\tau}, n)$.
  \item \emph{Per-fiber timeout}: a timeout on one fiber does not
    abort the entire verification; the result is recorded as
    ``inconclusive'' for that fiber, and the \v{C}ech computation
    accounts for it.
  \item \emph{Smaller search space}: each single-theory Z3 instance
    has a smaller search space than the combined-theory instance.
\end{enumerate}

\textbf{Objection 2: ``The number of fibers $k = 6^m$ is
exponential.''}  This is correct.  For $m = 3$ parameters with all
types relevant, $k = 216$.  However, duck-type inference
(Section~\ref{sec:opacity-classification}) typically reduces the
relevant fibers to 1--3 per parameter, giving $k \leq 27$ for
$m = 3$.  The implementation caps $k$ at~12 to prevent blowup.
Moreover, the $O(k)$ decomposition trades a constant-factor
increase in the number of queries for a qualitative simplification
of each query (single-theory vs.\ multi-theory).

\textbf{Objection 3: ``Lean doesn't use SMT; it uses its kernel.''}
Lean's kernel checks definitional equality by normalisation, not by
SMT.  However, for the \emph{specific} task of Python program
equivalence over $\PyObj$, Lean's kernel has no built-in theory
of $\PyObj$ arithmetic, so the user would need to either:
(a) implement the $\PyObj$ theory in Lean (substantial effort), or
(b) use Lean's SMT-integration tactic (\texttt{omega}, \texttt{decide})
which is limited to native Lean types and does not handle $\PyObj$.
In either case, the fiber decomposition advantage of CCTT applies
to whatever backend solver is ultimately used.

\textbf{What is genuinely new.}
The theorem establishes that \emph{type-fibered decomposition} of
verification problems is \emph{semantically sound} (via the \v{C}ech
complex and descent) and \emph{computationally beneficial} (by
avoiding theory combination).  The soundness guarantee — that local
verdicts compose to a global verdict iff $\Hone = 0$ — is not
available in traditional proof assistants, which lack a cohomological
framework for decomposing verification.
\end{remark}


%% ═══════════════════════════════════════════════════════════
\section{Theorem 3: The Normaliser Is a Decision Procedure for the
         Fold-Normal Fragment}
\label{sec:thm-normalizer}
%% ═══════════════════════════════════════════════════════════

The 7-phase normaliser implements a confluent and terminating rewrite
system on OTerms.  For a specific, well-characterised fragment of
programs, this rewrite system is \emph{complete}: two programs are
equivalent if and only if they normalise to the same canonical form.

\begin{definition}[Fold-normal fragment]
\label{def:fold-normal}
The \emph{fold-normal fragment} $\mathcal{F}_N$ consists of program
pairs $(f, g) \in \mathcal{L}_{\mathrm{OT}}^2$ whose normalised
OTerms $\mathrm{nf}(f), \mathrm{nf}(g)$ are built from the following
constructors only:
\begin{enumerate}[label=(\roman*)]
  \item $\mathsf{OVar}$, $\mathsf{OLit}$ (variables and constants),
  \item $\mathsf{OOp}(o, \vec{t})$ where $o$ is one of the
    \emph{canonicalised} operation names:
    $\mathsf{add}, \mathsf{sub}, \mathsf{mult}, \mathsf{floordiv},
    \mathsf{mod}, \mathsf{sorted}, \mathsf{reversed},
    \mathsf{len}, \mathsf{sum}$,
  \item $\mathsf{OFold}(o, t_{\mathrm{init}}, t_{\mathrm{coll}})$
    where $o$ is one of $\mathsf{iadd}, \mathsf{imult}, \mathsf{or},
    \mathsf{and}$,
  \item $\mathsf{OMap}(\lambda, t)$ (mapped collections),
  \item $\mathsf{OCase}(g, t, f)$ (conditionals),
  \item $\mathsf{OSeq}(\vec{t})$ (sequences).
\end{enumerate}
The fragment $\mathcal{F}_N$ excludes $\mathsf{OFix}$ (recursion),
$\mathsf{OUnknown}$, $\mathsf{OAbstract}$, and $\mathsf{OQuotient}$.
\end{definition}

\begin{theorem}[Normaliser completeness on the fold-normal fragment]
\label{thm:normalizer-complete}
The 7-phase normaliser is a decision procedure for the fold-normal
fragment $\mathcal{F}_N$:
\[
  \forall\, (f, g) \in \mathcal{F}_N.\;
  \mathrm{Equiv}(f, g)
  \iff
  \mathrm{nf}(f).\mathrm{canon}()
  =
  \mathrm{nf}(g).\mathrm{canon}().
\]
The normalisation runs in $O(n^2 \log n)$ time, where $n$ is the
combined size of the input OTerms.
\end{theorem}

\begin{proof}
\textbf{Termination.}
Each phase of the normaliser either:
\begin{enumerate}[label=(\alph*)]
  \item \emph{reduces} term size (Phases~1, 2: constant folding,
    identity elimination, involution cancellation), or
  \item \emph{lexicographically decreases} a complexity measure
    (Phases~3, 4: fold canonicalization absorbs outer operations
    into the fold, decreasing nesting depth), or
  \item \emph{stabilises} in a single pass (Phases~5, 6, 7: symbol
    unification, quotient normalisation, and case ordering are
    deterministic maps on the term structure).
\end{enumerate}
The outer loop runs at most 8 iterations (the implementation bound),
and each iteration terminates because each phase is a single
pass over the term tree.

\textbf{Confluence.}
On $\mathcal{F}_N$, the rewrite rules are orthogonal (no critical
pairs between rules from different phases):
\begin{itemize}
  \item Phase~2 rules are conditional on syntactic patterns
    ($x + 0 \to x$, $x \times 1 \to x$, etc.)\ and do not overlap
    with fold canonicalization (Phase~3) because the rules have
    disjoint left-hand side patterns.
  \item Phase~3 replaces named builtins with folds ($\mathsf{sum}(xs)
    \to \mathsf{OFold}(\mathsf{iadd}, 0, xs)$) — these are one-way
    rewrites with no overlap.
  \item Phase~4 absorption laws ($\mathsf{len}(\mathsf{list}(x))
    \to \mathsf{len}(x)$) reduce nesting depth monotonically.
  \item Phases~5--7 are deterministic and idempotent on
    $\mathcal{F}_N$ (no $\mathsf{OQuotient}$ to normalise, symbol
    unification is a fixed mapping).
\end{itemize}
By the Newman lemma, termination plus local confluence implies
global confluence.  Hence the normal form is unique.

\textbf{Soundness (if $\mathrm{nf}(f) = \mathrm{nf}(g)$, then
$\mathrm{Equiv}(f,g)$).}
Each rewrite rule preserves denotational semantics.  This is verified
for each rule individually:
\begin{itemize}
  \item Identity rules: $x + 0 = x$, $x \times 1 = x$ are
    arithmetic identities.
  \item Fold canonicalization: $\mathsf{sum}(xs) = \mathsf{fold}(+, 0, xs)$
    by the definition of $\mathsf{sum}$ in Python.
  \item Absorption: $\mathsf{sorted}(\mathsf{sorted}(x)) =
    \mathsf{sorted}(x)$ (idempotence of sorting).
  \item Map-fold fusion: $\mathsf{fold}(op, e, \mathsf{map}(f, xs))
    = \mathsf{fold}(op \circ f, e, xs)$ by the universal property
    of catamorphisms.
\end{itemize}
Hence identical normal forms imply semantic equivalence.

\textbf{Completeness (if $\mathrm{Equiv}(f,g)$, then
$\mathrm{nf}(f) = \mathrm{nf}(g)$).}
On $\mathcal{F}_N$, every program computes a function that can be
expressed as a composition of the canonicalised operations (item~(ii)
of Definition~\ref{def:fold-normal}).  The rewrite system normalises:
\begin{itemize}
  \item All syntactic variants of summation ($\mathsf{sum}$,
    for-loop with $\mathsf{+=}$, $\mathsf{functools.reduce}$ with
    $\mathsf{operator.add}$) to $\mathsf{OFold}(\mathsf{iadd}, 0,
    \cdot)$.
  \item All syntactic variants of sorting ($\mathsf{sorted}(\cdot)$,
    $\mathsf{.sort}()$) to $\mathsf{OOp}(\mathsf{sorted}, \cdot)$.
  \item All syntactic variants of list construction (comprehension,
    $\mathsf{map}$, for-loop with $\mathsf{.append}$) to
    $\mathsf{OMap}(\lambda, \cdot)$ with canonicalised bodies.
  \item All algebraic redundancies ($x + 0$, $\mathsf{sorted}(
    \mathsf{sorted}(x))$, $\mathsf{len}(\mathsf{list}(x))$) to
    their simplified forms.
\end{itemize}
Since $\mathcal{F}_N$ excludes $\mathsf{OFix}$ (no recursion) and
$\mathsf{OUnknown}$ (no opaque terms), the rewrite system exhausts
all sources of syntactic variation.  Two semantically equivalent
programs in $\mathcal{F}_N$ must therefore normalise to the same
canonical form.

\textbf{Time complexity.}
The term tree has $O(n)$ nodes.  Each normalisation phase is a
single bottom-up pass: $O(n)$ per phase.  The outer loop runs
$O(1)$ iterations (at most 8).  The canon() comparison at each
iteration is $O(n \log n)$ (string comparison of the canonical form,
which has length $O(n)$).  Hence total time is $O(n)$ for the
normalisation plus $O(n \log n)$ for the comparisons, giving
$O(n \log n)$.  We write $O(n^2 \log n)$ as a conservative bound
accounting for possible term growth during intermediate phases
(e.g., de~Bruijn renaming may traverse the full term for each
bound variable, contributing an $O(n)$ factor).
\end{proof}

\begin{remark}[Adversarial analysis of Theorem~\ref{thm:normalizer-complete}]
\label{rem:adv-normalizer}
\textbf{Objection 1: ``Lean's \texttt{ring} tactic does the same for
ring identities.''}  Lean's \texttt{ring} tactic normalises ring
expressions and is complete for the equational theory of commutative
rings.  However, \texttt{ring} operates on \emph{Lean terms}, not
Python programs, and does not handle folds, maps, sorting, or list
operations.  The fold-normal fragment $\mathcal{F}_N$ is
\emph{strictly larger} than the ring fragment: it includes
$\mathsf{OFold}$ (reducing arbitrary collections),
$\mathsf{OMap}$ (transforming collections), and
$\mathsf{sorted}$/$\mathsf{reversed}$ (ordering operations).

\textbf{Objection 2: ``The fold-normal fragment is too small to be
useful.''}  The fragment $\mathcal{F}_N$ covers:
\begin{itemize}
  \item Sum/product of a list (via fold canonicalisation),
  \item List comprehensions vs.\ explicit loops (via map unification),
  \item Sorted-then-processed vs.\ process-then-sort (via absorption),
  \item Redundant list/tuple wrappers (via identity elimination),
  \item Arithmetic simplification (via ring rules).
\end{itemize}
These cover a substantial fraction of the ``stylistic variation''
dimension in practical code.  The fragment does \emph{not} cover
recursive algorithms (Fibonacci, tree traversals) or programs with
opaque library calls — these require the path search or Z3 backend.

\textbf{Objection 3: ``Completeness relies on the rewrite system
being confluent, which isn't formally verified.''}  This is correct.
The confluence argument in the proof is a mathematical argument based
on the structure of the rewrite rules, not a machine-checked proof.
A fully rigorous treatment would require formalising the rewrite
system in a proof assistant and verifying the critical pair lemma.
We consider the argument to be at the level of rigour standard in
the term rewriting literature.
\end{remark}


%% ═══════════════════════════════════════════════════════════
\section{Non-Theorems: Advantages That Do Not Hold}
\label{sec:non-theorems}
%% ═══════════════════════════════════════════════════════════

Intellectual honesty requires acknowledging candidate advantages that
\emph{do not} hold up under adversarial scrutiny.

\subsection{Non-Theorem: CCTT has better asymptotic complexity than
Lean for all verification tasks}

This is false.  CCTT is a \emph{specialised} system for Python
program equivalence.  Lean is a \emph{general-purpose} proof
assistant.  For tasks outside the OTerm fragment — proofs about
arbitrary mathematical objects, formalisation of mathematical
theorems, or verification of algorithms in languages other than
Python — Lean is strictly more capable.  Even within the OTerm
fragment, Lean's soundness guarantee (small trusted kernel,
type-theoretic foundations) is stronger than CCTT's (which relies on
the correctness of Z3, the OTerm compiler, and the normaliser).

\subsection{Non-Theorem: CCTT's path search is asymptotically faster
than Lean's \texttt{aesop}}

\texttt{aesop} (Automated Extensible Search for Obvious Proofs) is
Lean~4's automated tactic for bounded proof search.  It performs a
best-first search over a user-configurable set of lemmas, similar in
spirit to CCTT's bidirectional BFS over the 24 dichotomy axioms.

Both systems have the same asymptotic structure:
\[
  T_{\mathrm{CCTT}} = O(|\mathcal{A}|^d \cdot n),
  \qquad
  T_{\mathrm{aesop}} = O(|\mathcal{L}|^d \cdot n),
\]
where $|\mathcal{A}|$ is the number of CCTT axioms, $|\mathcal{L}|$
is the number of \texttt{aesop} lemmas, and $d$ is the search depth.
The difference is in the \emph{axiom set}, not the \emph{algorithm}.
CCTT's axioms are curated for Python program equivalence; \texttt{aesop}'s
lemmas are general Lean library facts.  This makes CCTT more effective
\emph{on its domain} (fewer irrelevant axioms, less branching), but
the asymptotic complexity class is the same.

\subsection{Non-Theorem: Spec identification is a novel advantage}

CCTT's specification identification (D20) recognises that a term
computes ``factorial'' or ``Fibonacci'' and short-circuits verification.
This is effective but not novel — any system could implement
pattern matching on canonical forms.  The mathematical content is
trivial: if $f$ and $g$ both match a known pattern, and the pattern
uniquely determines a function, then $f = g$.  This is an engineering
optimisation, not a theoretical advantage.


%% ═══════════════════════════════════════════════════════════
\section{Summary of Comparative Advantages}
\label{sec:adv-summary}
%% ═══════════════════════════════════════════════════════════

\begin{center}
\renewcommand{\arraystretch}{1.3}
\begin{tabular}{>{\raggedright\arraybackslash}p{4cm}
                >{\raggedright\arraybackslash}p{4cm}
                >{\raggedright\arraybackslash}p{4.5cm}}
  \toprule
  \textbf{Advantage} & \textbf{CCTT} & \textbf{Lean / F${}^*$} \\
  \midrule
  Human annotation effort (Thm.~\ref{thm:zero-annotation})
    & $0$ tokens
    & $\Omega(n)$ tokens (translation) \\[3pt]
  Per-fiber theory isolation (Thm.~\ref{thm:fiber-decomp})
    & Single-sorted SMT per fiber
    & Combined-theory SMT (if SMT used) or no SMT \\[3pt]
  Decision procedure for $\mathcal{F}_N$ (Thm.~\ref{thm:normalizer-complete})
    & Built-in, $O(n^2 \log n)$
    & Not available; requires custom tactic \\[3pt]
  \midrule
  Soundness
    & Relies on Z3 + compiler + normaliser
    & Small trusted kernel \\[3pt]
  Generality
    & OTerm fragment of Python only
    & Arbitrary mathematics \\[3pt]
  Recursion handling
    & OFix (limited)
    & Full structural/well-founded recursion \\
  \bottomrule
\end{tabular}
\end{center}

The table makes explicit that CCTT's advantages are \emph{domain-specific}:
they hold for Python program equivalence on the OTerm fragment and
arise from specialisation (custom normaliser, curated axioms,
cohomological decomposition), not from a general superiority of the
underlying theory.  The price of specialisation is narrower scope and
a larger trusted computing base.

\begin{remark}[The specialisation--generality tradeoff]
The relationship between CCTT and Lean is analogous to the
relationship between a domain-specific language (DSL) and a
general-purpose language.  A DSL is more concise and more efficient
\emph{on its domain} but cannot express programs outside that domain.
CCTT is a ``DSL for Python equivalence verification'' — it automates
what general-purpose systems require human effort for, at the cost
of being unable to verify anything else.

This is a legitimate and well-understood tradeoff in formal methods.
The contribution of this chapter is to make the tradeoff
\emph{quantitative}: Theorems~\ref{thm:zero-annotation},
\ref{thm:fiber-decomp}, and \ref{thm:normalizer-complete} give
precise bounds on the effort and complexity savings that
specialisation provides.
\end{remark}
